\documentclass[12pt]{extarticle}

% ===== Encodage & langue =====
\usepackage{fontspec}
\usepackage{polyglossia}
\setdefaultlanguage{french}

% ===== Police Verdana (compiler avec XeLaTeX ou LuaLaTeX) =====
\setsansfont{Verdana}
\renewcommand{\familydefault}{\sfdefault}

% ===== Interligne =====
\usepackage{setspace}
\setstretch{1.1}

% ===== Marges réduites =====
\usepackage{geometry}
\geometry{a4paper, top=1.8cm, bottom=1.8cm, left=2cm, right=2cm}

% ===== Mathématiques =====
\usepackage{amsmath}

% ===== Espacement réduit autour des équations =====
\setlength{\abovedisplayskip}{4pt}
\setlength{\belowdisplayskip}{4pt}
\setlength{\abovedisplayshortskip}{2pt}
\setlength{\belowdisplayshortskip}{2pt}

% ===== Espacement réduit des sections =====
\usepackage{titlesec}
\titlespacing*{\section}{0pt}{1.2ex}{0.6ex}

% ===== Titre du document =====
\title{\textbf{Correction du contrôle séquence 1 version 4}}
\author{}
\date{}

\begin{document}
\maketitle

% --------------------------------------------------
\section*{Exercice 1}
\begin{align*}
19 - 5 - 3 + 21 - 6 &= 14 - 3 + 21 - 6 \\
                     &= 11 + 21 - 6 \\
                     &= 32 - 6 \\
                     &= 26
\end{align*}

% --------------------------------------------------
\section*{Exercice 2}

\begin{align*}
\textbf{a)}
6 + 4 \times 5 &= 6 + 20 \\
               &= 26
\end{align*}

\textbf{b)}
\begin{align*}
38 - 5 \times 4 &= 38 - 20 \\
                &= 18
\end{align*}

\textbf{c)}
\begin{align*}
31 + 16 \div 8 &= 31 + 2 \\
               &= 33
\end{align*}

\textbf{d)}
\begin{align*}
7 + 6 \times 4 + 5 - 9 \div 3 &= 7 + 24 + 5 - 3 \\
                              &= 33
\end{align*}

% --------------------------------------------------
\section*{Exercice 3}
\begin{align*}
4 \times 6 + 3 \times 7 &= 24 + 21 \\
                        &= 45
\end{align*}

La dépense est de 45~€.

% --------------------------------------------------
\section*{Exercice 4}

\textbf{a)} $29 - (4 + 8) \div 2$
\begin{align*}
&= 29 - 12 \div 2 \\
&= 29 - 6 \\
&= 23
\end{align*}

\textbf{b)} $(6 + 7) \times 5 - 8 + 4$
\begin{align*}
&= 13 \times 5 - 8 + 4 \\
&= 65 - 8 + 4 \\
&= 61
\end{align*}

\textbf{c)} $(25 - (9 + 6)) \times 4$
\begin{align*}
&= (25 - 15) \times 4 \\
&= 10 \times 4 \\
&= 40
\end{align*}

\textbf{d)} $5 \times (24 - (7 + 12)) - 6$
\begin{align*}
&= 5 \times (24 - 19) - 6 \\
&= 5 \times 5 - 6 \\
&= 25 - 6 \\
&= 19
\end{align*}

% --------------------------------------------------
\section*{Exercice 5}
\begin{align*}
13 - 4 \times 3 &= 13 - 12 \\
                &= 1
\end{align*}

Une gomme coûte 1~€.

% --------------------------------------------------
\section*{Exercice bonus 1}

\textbf{a)} $6 \times 5 + 14$

\textbf{b)} $8 \times (6 + 3)$

% --------------------------------------------------
\section*{Exercice bonus 2}

\textbf{a)} La somme du produit de 6 par 5 et de 9.

\textbf{b)} Le produit de la somme de 9 et de 2 par 7.

\end{document}